Estas son algunas preguntas que surgen de la práctica, con su respuesta.

\textbf{1. ¿Por qué la imagen de un espejo plano es virtual?}

Porque los rayos reflejados nunca pasan por el punto donde vemos la imagen. Solo sus prolongaciones hacia atrás del espejo se cruzan ahí, y por eso no se puede proyectar sobre una pantalla.

\textbf{2. ¿Qué tamaño mínimo debe tener un espejo para que una persona se vea completa?}

La mitad de su estatura. El rayo que va de los pies a los ojos rebota en un punto que está justo a la mitad de esa altura, y lo mismo pasa con el de la cabeza, sin importar a qué distancia esté la persona del espejo.

\textbf{3. ¿El espejo invierte la izquierda y la derecha?}

No exactamente. Lo que invierte es el adelante y el atrás, en la dirección perpendicular al espejo. Si levantamos la mano derecha, la imagen levanta la mano que está del mismo lado, pero como la imagen nos mira de frente, nos parece su mano izquierda.

\textbf{4. ¿Cuántas imágenes se forman con dos espejos planos que forman un ángulo?}

Si el ángulo entre ellos es $\alpha$ y $360/\alpha$ es un número entero, se forman
\begin{equation}
N = \frac{360}{\alpha} - 1
\label{eq:imagenes}
\end{equation}
imágenes. Con dos espejos a 90 grados, por ejemplo, se ven tres.
